\section{Directed-cycle consequences}\label{sec:cycles}

\begin{contract}
\textbf{Assumes:} \cref{thm:main}. \quad
\textbf{Produces:} degree and order bounds for oriented graphs without a directed triangle, including
a minimum-degree threshold of $n/3$ for the existence of a directed triangle
(\cref{cor:cycles}). \textbf{Why:} these are the standard corollaries; they show what the theorem buys
and they are the reason the conjecture is studied alongside the Caccetta--H\"aggkvist problem.
Nothing here is used earlier.
\end{contract}

The mechanism is one observation: in an oriented graph without a directed triangle, no second
out-neighbor can be an in-neighbor, because a path $v\to u\to w$ together with the arc $w\to v$ would
be a directed triangle. So the second out-neighborhood sits inside the set of \emph{non-neighbors},
and \cref{thm:main} becomes a counting statement about three disjoint classes of vertices.

\begin{corollary}[Non-neighbors and directed cycles]\label{cor:cycles}
Let $D$ be a nonempty finite oriented graph of order $n$. For $v\in V(D)$ put
\[
 N_1^-(v)=\{u:u\to v\},\qquad d^\pm(v)=\abs{N_1^\pm(v)},\qquad
 \delta^\pm(D)=\min_{x\in V(D)}d^\pm(x),
\]
and let
\[
 \overline N(v)=V(D)\setminus\bigl(\{v\}\cup N_1^+(v)\cup N_1^-(v)\bigr)
\]
be the set of \emph{non-neighbors} of $v$ (vertices adjacent to $v$ in neither direction). Then:
\begin{enumerate}[label=\textup{(\roman*)},leftmargin=1.6em,itemsep=3pt]
\item if $D$ has no directed $3$-cycle, then some vertex $v$ satisfies
  \[
   d^+(v)\le\abs{\overline N(v)},\qquad n-1\ge2d^+(v)+d^-(v);
  \]
\item if $\delta^+(D)\ge n/3$ and $\delta^-(D)\ge n/3$, then $D$ contains a directed $3$-cycle;
\item if $r\ge1$ is an integer, $d^+(v)=d^-(v)=r$ for every vertex, and $D$ has directed girth
  exactly four, then $n\ge3r+1$.
\end{enumerate}
Here the directed girth is the length of a shortest directed cycle.
\end{corollary}

\inwords{Without a directed triangle, some vertex has at most as many out-neighbors as non-neighbors,
which forces $n-1\ge2d^++d^-$; hence minimum in- and out-degree at least $n/3$ forces a directed
triangle, and a regular oriented graph of girth four satisfies $n\ge3r+1$.}

\usedin{Not used in the proof of \cref{thm:main}; these are consequences, and \cref{sec:history}
records where they appear in the literature.}

\begin{proof}
Assume that $D$ has no directed $3$-cycle, and let $w\in N_2^+(v)$; choose a path $v\to u\to w$. By
the definition of exact distance, $w\ne v$ and $w\notin N_1^+(v)$. If also $w\in N_1^-(v)$, then
$w\to v$, and $v\to u\to w\to v$ is a directed $3$-cycle --- impossible. Hence
\begin{equation}\label{eq:N2-in-Z}
 N_2^+(v)\subseteq\overline N(v)\qquad\text{for every }v .
\end{equation}
By \cref{thm:main} there is a vertex $v$ with
\[
 d^+(v)=\abs{N_1^+(v)}\le\abs{N_2^+(v)}\overset{\eqref{eq:N2-in-Z}}{\le}\abs{\overline N(v)} .
\]
Since $D$ is oriented, $N_1^+(v)$ and $N_1^-(v)$ are disjoint and both avoid $v$, so the three sets in
the definition of $\overline N(v)$ partition $V(D)\setminus\{v\}$:
\begin{equation}\label{eq:partition}
 n-1=d^+(v)+d^-(v)+\abs{\overline N(v)}\ge 2d^+(v)+d^-(v),
\end{equation}
which is (i); \cref{fig:cycles} draws the partition.

For (ii), if $D$ had no directed $3$-cycle then (i) would give
\[
 n-1\ge2\delta^+(D)+\delta^-(D)\ge2\cdot\frac n3+\frac n3=n,
\]
a contradiction. (The degrees are integers, so $\delta^\pm(D)\ge n/3$ means
$\delta^\pm(D)\ge\lceil n/3\rceil$; the argument covers the endpoint when $3$ divides $n$.)

For (iii), directed girth four excludes directed triangles, and $d^+=d^-=r$ turns
\eqref{eq:partition} into $n-1\ge3r$.
\end{proof}

\begin{figure}[t]
\centering
\begin{tikzpicture}[
  v/.style={circle,draw=black!70,fill=black!8,inner sep=0pt,minimum size=6mm,font=\scriptsize},
  o/.style={circle,draw=blue!55!black,fill=blue!18,inner sep=0pt,minimum size=5.2mm,font=\scriptsize},
  i/.style={circle,draw=orange!60!black,fill=orange!20,inner sep=0pt,minimum size=5.2mm,font=\scriptsize},
  nn/.style={circle,draw=green!40!black,fill=green!12,inner sep=0pt,minimum size=5.2mm,font=\scriptsize},
  ar/.style={-{Stealth[length=1.4mm]},draw=black!70},
  lab/.style={font=\scriptsize,text=black!75}
]
% regions
\draw[green!45!black,dashed,rounded corners=6pt] (-1.7,2.1) rectangle (3.4,3.8);
\draw[orange!65!black,dashed,rounded corners=6pt] (-2.1,-1.5) rectangle (-0.9,1.5);
\draw[blue!55!black,dashed,rounded corners=6pt] (0.9,-1.5) rectangle (2.1,1.5);
% vertices
\node[v] (v) at (0,0) {$v$};
\node[o] (a) at (1.5,0.9) {};
\node[o] (b) at (1.5,-0.9) {};
\node[i] (c) at (-1.5,0.9) {};
\node[i] (d) at (-1.5,-0.9) {};
\node[nn] (e) at (0.1,2.9) {};
\node[nn] (f) at (1.5,2.5) {};
\node[nn] (g) at (2.8,3.0) {};
\node[nn] (h) at (-1.1,2.6) {};
% arcs
\draw[ar] (v) -- (a);
\draw[ar] (v) -- (b);
\draw[ar] (c) -- (v);
\draw[ar] (d) -- (v);
\draw[ar] (a) to[bend left=15] (f);
\draw[ar] (b) to[bend right=12] (g);
% labels
\node[lab,anchor=west,text=blue!50!black] at (2.25,0) {$N_1^+(v)$};
\node[lab,anchor=east,text=orange!55!black] at (-2.25,0) {$N_1^-(v)$};
\node[lab,anchor=south,text=green!35!black] at (0.85,3.85) {$\overline N(v)$: the non-neighbors};
\node[lab,anchor=north,text=green!35!black] at (1.6,2.05) {$N_2^+(v)\subseteq\overline N(v)$};
% the count
\node[lab,anchor=north] at (0,-1.85) {$n-1=d^+(v)+d^-(v)+\abs{\overline N(v)}$};
\node[lab,anchor=north] at (0,-2.5) {$\overset{\text{(i)}}{\ge}2d^+(v)+d^-(v)$};
\end{tikzpicture}

\caption{The partition behind \eqref{eq:partition}, in an oriented graph without a directed triangle.
\emph{Look at} the three classes around $v$: out-neighbors, in-neighbors and non-neighbors. The
label inside the third class records \eqref{eq:N2-in-Z}, that the second out-neighborhood
sits inside it: a second out-neighbor that was also an in-neighbor would close a directed
triangle. \emph{Why it matters:} \cref{thm:main} bounds the first class by the second
out-neighborhood, hence by the third class, and the three classes partition $V\setminus\{v\}$; this
gives $n-1\ge2d^++d^-$ with no further input.}
\label{fig:cycles}
\end{figure}

\begin{remark}[Where these statements come from]\label{rem:cycles-history}
Part (i) is the oriented-graph form of Thomass\'e's non-neighbor statement, recorded as
Conjecture~6.12 in Sullivan's survey~\cite{Sullivan2006}; the consequence for minimum in- and
out-degree is noted in Section~3 of the same survey, and the regular girth-four bound in Section~4,
where it is attributed to the Behzad--Chartrand--Wall minimum-order
conjecture~\cite{BehzadChartrandWall1970}. The point of recording them here is that they follow from
\cref{thm:main} alone, with no additional hypothesis on $D$.
\end{remark}
